\paragraph{Bounding the unscaled depth by covariant differential operators.}
\label{par:unscaled-light-depth}
All orders in this paragraph are in units of \(v_p(p)=1\).
They give bounds for \(\nu_W(\lambda)\) by the
\hyperref[par:slice-finite-depth]{definition of chart depth}.
A covariant differential operator on an integral germ may be evaluated over a complete algebraically closed valued extension of the chart field. Its value depends on a finite jet of the germ. Over the extension,
these jets are integral linear combinations of coefficient
monomials. The integral harmonic coefficient lattice
extends as well. Indeed, in the fixed multidegree the
harmonicity equations are linear over the original
valuation ring. Their kernel extends under flat scalar
extension, since the new valuation ring is torsion-free
over the original DVR.

If \(H\) is the matrix of the harmonicity equations and \(V\subset V'\) the extension of valuation rings, this says
\[
  (\ker H)\otimes_V V'=\ker(H\otimes_V V').
\]
 The covariance operators commute
with coefficient extension. By bilinearity and the
ultrametric inequality, each lower bound on the original valuations also holds after scalar extension, in the same units. Thus any extended test giving a finite valuation bounds the original infimum from above.
\label{cor:slice-unscaled-depth-1}

\begin{proposition}[Upper bounds for the unscaled chart depth]
\label{prop:primitive-light-covariants}
At every allowed finite place, primitive columns in the local
eigenspaces bound the unscaled depth \(\nu_W(\lambda)\) used in
Proposition~\ref{prop:slice-finite-readings}. Use the
\hyperref[par:slice-finite-depth]{weighted chart depth} with \(v_p(p)=1\).
At the small primes the eigenflag and colength estimate is
\eqref{eq:light-small-prime-depth}; at \(q,d\) the integral
Hodge--Tate estimate is \(\nu_W\le D_L/(p-1)\). On the good charts
the bounds use monomials or alternating pairs of Frobenius orbits. In the balanced level-one conjugating case write \(C_{\rm lead}\) for the leading flow cost. The degree contributions \(D_a\) enter as
\[
 C_{\rm lead}=\frac1{p^2-1}\sum_a D_a.
\]

\end{proposition}

\begin{proof}
\proofstep{Choice of primitive columns in each eigenspace.}
For each label \(\tau\) of rank \(r_\tau>0\), let
\(a_{\tau,1}\ge\cdots\ge a_{\tau,r_\tau}\ge0\) be the
light row weights, and put \(a_{\tau,r_\tau+1}=0\).
Choose \(r_\tau\) columns of integral characters on that
frame, with disjoint column sets for distinct labels.
The chosen number of rows permits this. If their derivative
matrix has leading minors \(\Delta_{\tau,i}\), use the
product of the corresponding determinant polynomials to
the powers \(a_{\tau,i}-a_{\tau,i+1}\).
Row \(j\) occurs precisely in the minors of sizes \(i\ge j\).
Its exponent is therefore
\[
 \sum_{i=j}^{r_\tau}(a_{\tau,i}-a_{\tau,i+1})=a_{\tau,j}.
\] A harmonic
contraction pairs conjugate types in the same column.
Each column belongs to a single type, so every contraction
kills the polynomial. Its coefficients are integral.

In formal logarithm coordinates the character derivatives
evaluate the determinant product with no factorial loss.
If \(v_p(\Delta_{\tau,i})\le h_\tau\) for each \(i\), this
product has valuation at most 
\[
  \sum_\tau a_{\tau,1}h_\tau.
\]
The integral covariant character-germ comparison replaces the
ordinary reading by this product plus an error of strictly
higher valuation. It uses the integral pairing and chooses
auxiliary eigenvalues congruent to one to accuracy greater
than the order of the nonzero product. Thus the covariant
reading has the same valuation.
\label{cor:slice-unscaled-depth-2}

For the bounds on the minors, start with a finite collection
of differential columns attaining the chart colength.
In the fixed first flag row choose an entry of least valuation
and use its column to clear that row in every other column
by integral multiples.
For a pivot \(a\ne0\) of least valuation and another entry \(b\) on its row, the clearing step uses
\[
  v_p(b/a)\ge0,\qquad b-(b/a)a=0 .
\]
 At each step take the minimum in the
next fixed flag row of the residual columns. Each residual pivot comes from an original
column after earlier column subtractions.

For the nested
sequence the leading determinants are unchanged by these
subtractions and have valuations equal to the partial sums
of pivot valuations. Since each pivot is integral, the
partial sum is at most the full sum, which is the valuation
of the determinant ideal.
If \(p_1,\ldots,p_n\) are the successive nonzero integral pivots, the cost of any initial segment is bounded by
\[
  \sum_{j=1}^{i}v_p(p_j)\le\sum_{j=1}^{n}v_p(p_j)\qquad(i\le n).
\]
 This choice of integral characters over the original ring gives the required leading minors.
\label{cor:slice-unscaled-depth-3}

\proofstep{The Hodge and colength inputs at the small primes.}
At the designated small primes use the split auxiliary
quadratic field. The selected CM half is partitioned by the
host places of the split blocks. If a block has \(n_w\)
embeddings and discriminant coefficient \(d_w\), the CM
character comparison on the integral chart bounds the
product of its absolute determinant lower bounds below by
\(|p|^{n_wq_j(1/(p-1)+d_w/2)}\).
Taking logarithms and dividing by \(\log|p|<0\), then using
\(\sum_w n_w=f\) and \(\sum_w n_w d_w=fv_d(p)\), gives
\[
        \sum_{\tau\text{ in the selected half}} h_\tau
            \le fq_j(1+v_d(p)/2).
\]
On the complementary half the integral Hodge--Tate
character inclusion contains \(p\) times the tangent
dual, since \(v_p(p)=1\ge1/(p-1)\).

Its restriction projects onto \(p\) times the dual of the
integral rank-\(q_j\) eigendirection frame, with determinant
cost at most \(q_j\) for each label.
The small-prime page bounds the first light row of the
selected half by \((2b+S_c)/T\), and the sum of the first
rows of the complementary half by \(f\eta_0/T\).
With \(S_e=S_c+\eta_0\), \(q_j\le m\) and
\(v_d(p)\le v_d\), the determinant product gives
\begin{equation}\label{eq:light-small-prime-depth}
              T\nu_W\le m(2+v_d)(2fb+fS_e).
\end{equation}
This uses the summed colength once on the selected half.
The global comparison may use \(2fb\) as a majorant for
the weight sum of any lift. At \(q,d\) the crude integral
Hodge--Tate bound, obtained by using the degree of each
label in the determinant reading, is

\[
  \nu_W\le D_L/(p-1).
\]
The CM colength and Hodge--Tate inclusions are local inputs
on the given uniformization and eigendirection frames.
\label{cor:slice-unscaled-depth-4}

\paragraph{Depth bounds at good primes and odd shifts.}
On the unramified chart, order the embedding labels by their
residue Frobenius. Let \(C_\tau\) be the character colength at label \(\tau\). If \(\tau_h\) denotes its \(h\)-th successor, the unramified estimate uses the ranks of the conjugate eigenspaces
\[
  C_\tau\le\sum_{h\ge1}p^{-h}r_{\overline{\tau_h}}.
\]
 The row systems use this same
labelled orbit. The characteristic-zero frame and the
integral chart give the character and monomial blocks with
the prescribed types.
\label{cor:slice-unscaled-depth-5}

\proofstep{The diagonal monomial on a split block.}
For a monomial block use formal curve coordinates with
identity linear part on the rows of the indicated half.
The product of the leading minor polynomials has
diagonal monomial with coefficient one. For each minor,
the identity permutation uniquely maximizes the sum of
the products of the row and column indices, by the
rearrangement inequality.
For increasing row and column sequences \((r_i),(c_i)\), strictness for a nonidentity permutation reads
\[
  \sum_i r_ic_{\sigma(i)}<\sum_i r_ic_i\qquad(\sigma\ne1).
\]
 This identifies its diagonal
term uniquely even inside the product. If \(b_i\) are the diagonal monomial exponents, the
prescribed germ derivative reads \(\prod_i b_i!\) because
the coefficient of this monomial is one and the higher
total-degree terms give zero. The single-half covariant
factor is one. Legendre's formula gives the loss
\[
         \sum_i v_p(b_i!)=\sum_i\sum_{a\ge1}\lfloor b_i/p^a\rfloor
                 \le \frac{\sum_i b_i}{p-1},
\]
with zero loss when all \(b_i<p\).

If a first-row
exponent is at most \(s_0(p-1)\), its integral weight gaps
can be distributed among \(s_0\) row slots of capacity
\(p-1\), by filling the slots successively and splitting
gaps at the boundaries. Later rows use only a subset of
the gaps. Different labels of one type reuse the slots
using their labelled coordinates. When 
\[
  4rf s_0\le n_W,
\]
there are enough rows of each type to obtain a monomial of valuation zero on the whole chart. In particular, \(p>D_L\) permits \(s_0=1\).
In the split case the character block of exponent \(u\)
costs \(u m/(p-1)\) on the distinguished part and the
complementary monomials use the remaining degree.
\label{cor:slice-unscaled-depth-6}

For an alternating pair of orbits the selected inert-prime
hypothesis gives \(p\ge5\). Put \(n=k-1\ge1\) and let
\(d_0\) be its orbit length. An odd shift \(1\le i<d_0\) uses block
exponents \(p^i y,y\).
On that pair the total degree consumed is
\[
  p^iy+y=(p^i+1)y .
\]
 The adapted chain begins with at least
\(n-((i-1)/2+2)(m-j)\) independent residue directions,
loses at most \(m-j\) at each of \((i-1)/2\) successive
semilinear steps and loses at most \(j\) in the final
pairing. Hence at least \(n-(i+1)m\) rows remain.

An integral completion of the opposite-side vectors puts
their retained pairing in negative identity form. The
integral gauge estimates and the odd-level jet theorem on
this chart then give admissible \(y\) in every window of
the specified length, as well as \(y=0\). For positive
\(y\), write \(\gamma=p/(p^2-1)\). The block cost is at most
\[
  n\gamma(p^i+1)y-ny+n(1+\log_p n)
                 +\frac{(i+1)m}{p-1}(p^i+1)y .
\]
The term with \(m\) bounds the factorials of the lost
rows. At \(y=0\) the cost is zero. This uses the odd-level
jet estimate at the stated row ranks, normal form of its
mixed gauge coefficients and integral curve inputs.
\label{cor:slice-unscaled-depth-7}

\proofstep{The cost of the alternating orbit.}
To see how the blocks give the orbit price, write \(W_a\)
for the distinguished exponents at the light weight and
\(w_a>0\) for the target exponents. The degree comparison
gives \((1-\theta)w_a\le T W_a\le w_a+C_1S_e\) with
\(0\le\theta\le1\) and \(n\theta\le5m\).
Set \(V_a=(1-\theta)w_a/T\).
The directed odd shifts allocate nonnegative variables
\(x_{a,i}\) with constraints
\[
         \sum_i p^i x_{a,i}+\sum_i x_{a-i,i}\le nV_a .
\]

On demands \(D_a=nV_a\), define the cost
\(F_p(D)=\gamma\sum_a D_a-\max\sum_{a,i}x_{a,i}\).
Summing the constraints proves that the maximum is at
most \(\sum_a D_a/(p+1)\). The cost is nonnegative,
homogeneous and invariant under cyclic shifts.
If levels \(i\ge I=i_0(k,p)\) are discarded from a
maximizer, the lost flow is at most \(p^{-I}\sum_aD_a\).
Round each remaining \(x_{a,i}/n\) down to an allowed
integer \(y_{a,i}\), with loss at most the window size
\(S\) in this quotient. Then
\[
 u_a=W_a-\sum_i p^i y_{a,i}-\sum_i y_{a-i,i}\ge0
\]
is an integer character-block exponent.

In the total cost
the leading character term and leading odd terms sum to
\(n\gamma\sum_a W_a\), by the displayed equality and
cyclic reindexing. The saving from the rounded flow loses
at most \(d_0 I n S\). The positive odd factorial costs
are at most \(2I m\sum_a V_a/(p-1)\) since
\(p^i+1\le2p^i\) and \(i+1\le I\) on the retained levels.
The character cost beyond its leading term is at most
\(m\sum_aW_a/(p-1)\) and the additive odd cost is at most
\(d_0 I n(1+\log_p n)\).
The chart range with these odd blocks has
\[
 I\le5\log k,\qquad S+1+\log_p n\le707 n^2\log k.
\]
Multiplying the rounding and additive window costs by \(T\) gives the combined upper bound
\[
     3535 d_0 Tn^3(\log k)^2 .
\]

The discarded levels cost at most
\(p^{-I}k\sum_a w_a\).
Replacing \(V_a\) by \(W_a\) and paying the remaining
monomial degrees adds at most a fixed multiple of
\((m(\sum_a w_a)\log k+d_0 kS_e)/p\) on the light-weight
and row-budget comparisons. For the sum over orbits, the total
length is at most \(f\). Use the uniform weight bound
\(\widehat W=2fb\). Below the threshold
\(p_0=mb/(Tf n^3\log b)\),
\[
 Tf n^3(\log k)^2
 < \frac{mb(\log k)^2}{p\log b}
 \le \frac{m\widehat W\log k}{p}.
\]
Thus the absolute window loss fits the weight allowance with
its fixed numerical multiplier. At \(p\ge p_0\) the typed row
supply on the selected chart gives the zero-depth monomial
allocation, including equality at the threshold.
\label{cor:slice-unscaled-depth-8}

\proofstep{Vanishing of the depth at large nonsplit primes.}
At large nonsplit primes the leading cost has a direct
bound. If \(d_0\ge2\) and all demand ratios are at most \(p\), the system
\(p x_a+x_{a-1}=D_a\) has the cyclic solution
\[
 x_a=\frac{\sum_{i=0}^{d_0-1}(-p^{-1})^iD_{a-i}}
                {p(1-(-p^{-1})^{d_0})}.
\]
Pair successive terms in the numerator. Each pair is
nonnegative and in odd length the leftover term is
positive for positive demands. This realizes the upper
bound on the total flow using only level one. Thus
\[
 F_p(D)=\sum_a D_a/(p^2-1).
\]
The selected target ratios give the required balance for
\(p>X_2\), so the leading cost is at most
\(2k\sum_a w_a/p^2\).

At smaller primes the forced content is the orbit average
of the orbit programmes. Cyclic invariance gives their sum
over representatives in the place formula. When a
Frobenius power is conjugation, choose the first
admissible odd level on the same typed chart and the
programme comparison for that orbit. For Frobenius equal to
conjugation choose level one directly and write

\[
  W=(p+1)y+u
\]
 with \(y\) in the jet window and
\(0\le u\le(p+1)S\). This gives the explicit
\(1/(p^2-1)\) leading term in that case.
The cost comparison for the first conjugating odd level
and the block supplies in the respective charts are used
with these hypotheses.
\label{cor:slice-unscaled-depth-9}
\end{proof}
